\chapter{Probabilité}
\section{Vocabulaire : Expérience aléatoire,issue,univers,...}
\lipsum[2]
\subsection{Définition : Expérience aléatoire}

Une expérience est dite aléatoire lorsque l'on ne peut pas prévoir l'issue de cette expérience.

\subsection{Vocabulaire}

plop\sidenote{ds sd  dddd} plop

\lipsum[3]
\begin{itemize}
	\item \ul{Issue} : Un résultat de l'expérience aléatoire
	\item $50$ ; $1000$ ou ${Jackpot}$ ...
	\item Univers : L'ensemble de toutes les issues possibles
	      \begin{enumerate}
		      \item $\Omega=\left\{{LOSE} ; 50 ; 100 ;... ;1000 ;{JACKPOT}\right\}$
		      \item **Événement** : Un ensemble d'issues
		      \item $A$ : "Obtenir - de 450\eur" $\Rightarrow A=\left\{{LOSE} ;50 ;100 ;200 ;300\right\}$
		      \item **Expérience aléatoire** : "Choisir,au hasard,une lettre dans l'alphabet"
	      \end{enumerate}
	\item $\Omega =\left\{A;B;C;D;E;F;G;H;...;X;Y;Z\right\}$
	\item $E_1=\left\{A;E;I;O;U\right\}$ est un **événement**.
	\item En français, cet événement se traduit par : $E_1: {"La lettre choisie est une voyelle"}$
	\item $E_2=\left\{K;W;X;Y;Z\right\}$ est un autre **événement**.
	\item Ce second événement se traduit par : $E_2: {"La lettre choisie vaut 10 pts"}$
\end{itemize}

